\documentclass[11pt]{article}
\usepackage[margin=1in]{geometry}
\usepackage{amsmath,amssymb,amsthm,booktabs}
\usepackage[hidelinks]{hyperref}
\usepackage{xurl}
\setlength{\emergencystretch}{3em}
\newtheorem{theorem}{Theorem}
\newtheorem{lemma}[theorem]{Lemma}
\newtheorem{corollary}[theorem]{Corollary}
\newcommand{\mr}{\operatorname{mr}}
\newcommand{\ar}{\operatorname{ar}}
\newcommand{\rk}{\operatorname{rank}}
\newcommand{\C}{\mathbb C}
\newcommand{\F}{\mathbb F}
\title{Prime four-bridge knots with augmentation rank three}
\author{}
\date{17 November 2026}
\begin{document}
\maketitle
\begin{abstract}
We give an explicit twist family of prime knots with meridional rank and
bridge number four, complex augmentation rank three, and crossing number
$17+6n$ for every integer $n\ge0$. Their genera are $7+3n$.
A full alternating-group quotient survives the twists, while a local
rank-one obstruction survives because Hurwitz moves preserve the ordered
product. An explicit three-dimensional finite reflection image attains
the augmentation bound. A nondegenerate extension of a possibly singular
Gram form also bounds every Coxeter reflection-class quotient by three
generators. Thus neither one such Coxeter quotient nor one finite complex
rank-one reflection image gives the sharp meridional lower bound for
these prime knots. Connected sums with trefoils retain a gap of one.
\end{abstract}

Let $G_K=\pi_1(S^3\setminus K)$, let $\mr(K)$ denote the minimum number
of meridians generating $G_K$, and let $b(K)$ be bridge number. Ordinary
complex augmentation rank is denoted $\ar(K,\C)$; we use field-valued
$U=1$ cord augmentations with meridian value different from one and
arbitrary nonzero longitude value, as in Cornwell
\cite[Definition~1.8]{cornwell}.
For an integer $n\ge0$, define $K_n$ to be the closure of
\begin{equation}\label{eq:braid}
 \beta_n=\sigma_1\sigma_2\sigma_3^{-2}\sigma_2
 \sigma_3^{-2}\sigma_2\sigma_1^2\sigma_2^3
 \sigma_3^{-1}\sigma_2^{3+6n}\ \in B_4.
\end{equation}

\begin{theorem}\label{thm:main}
For every integer $n\ge0$, $K_n$ is a prime knot and
\[
 \mr(K_n)=b(K_n)=4,\qquad g(K_n)=7+3n,\qquad c(K_n)=17+6n,
 \qquad \ar(K_n,\C)=3.
\]
The largest meridian-class generating rank supplied by a finite complex
image whose meridian displacement has rank one is exactly three.
Every Coxeter quotient taking meridians to reflections has Coxeter rank
and meridian-class generating rank at most three; the Coxeter group need
not be finite.
\end{theorem}

The obstruction behind the theorem compares the root span and the
covector span independently. Four independent roots and four independent
covectors would force two local folds, leaving a fixed three-direction
system known to be inconsistent. The apparent dependence on the last
three crossings disappears: the remaining tail preserves an ordered
product, so the same contradiction holds for every tail in
\eqref{eq:braid}. The twist length is chosen to preserve a separate full
$A_8$ quotient and a positive three-dimensional representation.
Primeness follows from the homogeneous braid theorem of Feller,
Lewark and Orbegozo Rodriguez \cite{flo}, using two run counts which
do not change under these twists.

Strict augmentation gaps themselves precede this construction:
Cornwell's Corollary~1.12 gives augmentation rank two and meridional rank
three for several knots \cite{cornwell}. The contribution here is this
explicit prime family, its uniform paired-span obstruction, and the
extension of that obstruction to singular canonical Coxeter forms.
The finite complex comparison theorem \cite[Theorem
\texttt{thm:domination}]{domination} supplies the finite-class consequence;
it is not a new classification theorem about finite reflection groups.
We do not assert hyperbolicity or a non-graph double branched cover for
every member of the family.

\section{The alternating-group quotient and four bridges}

Use the literal strand convention
\begin{equation}\label{eq:hurwitz}
 T(a,b)=(b,b^{-1}ab)\quad\text{for }\sigma_i,
 \qquad H(a,b)=(aba^{-1},a)\quad\text{for }\sigma_i^{-1}.
\end{equation}
These are inverse maps and preserve the ordered product $ab$.
For involutions, putting $q=ab$ gives
\[
 H^k(a,b)=(q^ka,q^{k-1}a).
\]
This follows for $k=0$ and is preserved by one substitution. In
particular $q^d=1$ implies $H^d=T^d=1$ on this pair.

Place the following permutations on the four top strands, with
permutations composed from right to left:
\begin{equation}\label{eq:a8}
 x_0=(01)(23),\quad x_1=(14)(23),\quad
 x_2=(05)(16),\quad x_3=(12)(67).
\end{equation}
They generate $A_8$. One exact certificate is breadth-first enumeration
of their products: it gives $20160=|A_8|$ distinct even permutations.
This is a finite exhaustive computation on explicitly specified
permutations, reproduced by the accompanying verifier.
Let $U$ be the first fourteen letters of $\beta_n$. At the end of $U$
the middle pair in this action is
\[
 a=(05)(14),\qquad b=(16)(23),\qquad
 ab=(05)(164)(23),
\]
whose product has order six. The four labels return after $U\sigma_2^3$;
this too is a direct substitution in \eqref{eq:hurwitz}.
The period-six identity proves return for $U\sigma_2^{3+6n}$ for every
$n\ge0$. Thus \eqref{eq:a8} defines a surjection
$G_{K_n}\twoheadrightarrow A_8$ taking meridians into the class $C$
of double transpositions.

Three double transpositions cannot generate a transitive group on eight
points. Indeed, their six transposition factors provide at most six
edges in a graph on eight vertices, so the graph is disconnected and
its components are preserved. The four displayed permutations do
generate $A_8$; consequently $d_C(A_8)=4$, and $\mr(K_n)\ge4$.
The permutation of $\beta_0$ is a four-cycle. Adding six adjacent
transpositions leaves it unchanged, so every closure is connected.
The elementary closed-braid bridge bound gives
$\mr(K_n)\le b(K_n)\le4$, proving equality. In particular the braid
index is also four.

\section{Primeness and exact size}

All appearances of $\sigma_1,\sigma_2$ in \eqref{eq:braid} are positive,
and all appearances of $\sigma_3$ are negative, so the braid word is
homogeneous. For $i=2,3$, delete all letters except those indexed $i-1$
or $i$ and count the alternating nonempty runs of equal generator
indices, linearly in the written word. The restricted words have run
patterns
\[
 1\mid2^3\mid1^2\mid2^{6+6n},\qquad
 2\mid3^2\mid2\mid3^2\mid2^4\mid3\mid2^{3+6n},
\]
respectively, with signs suppressed. Hence $g_2=4$ and $g_3=7$ for
all $n$. Lemma~5.1(i) of \cite[printed p.~20]{flo} says that the braid
diagram admits a decomposition circle precisely when some such count
is two or three. There is no decomposition circle here. Theorem~1.1 of
\cite[printed p.~1]{flo} states that a nonprime homogeneous braid closure
has such a circle in this diagram. Thus $K_n$ is prime. It is nontrivial
because its bridge number is four. These are linear run counts; joining
the beginning and end of a restricted word would give the wrong test.

We record a uniform Alexander identity, both to distinguish the knots
and to determine their genus without a genus-minimality import for
homogeneous diagrams. Write
\begin{align*}
 P(t)&=1-3t+4t^2-8t^4+18t^5-21t^6+15t^7-6t^8+t^9,\\
 P^*(t)&=t^9P(t^{-1}).
\end{align*}
\begin{lemma}\label{lem:alexander}
The reciprocal normalization with constant coefficient one is
\begin{equation}\label{eq:alexander}
 \Delta_{K_n}(t)=\frac{P(t)+t^{6n+6}P^*(t)}{1+t}.
\end{equation}
Its degree is $14+6n$, its value at one is one, and its determinant is
$117+366n$.
\end{lemma}
\begin{proof}
We give the symbolic calculation underlying all parameter values.
Let $G_i(t)$ be the identity matrix of order four with its $i,i+1$ block
replaced by
\[
 \begin{pmatrix}1-t&t\\1&0\end{pmatrix},
\]
and form $B_U$ by multiplying the fourteen matrices for $U$ from left
to right, using $G_i^{-1}$ for negative letters.
Put $N=G_2-I$. Direct multiplication gives
\[
 N^2=-(1+t)N,\qquad
 G_2^k=I+S_kN,\qquad S_k=\sum_{j=0}^{k-1}(-t)^j.
\]
The second identity follows by induction. Since $N$ has rank one, the
cofactor obtained by deleting the last row and column of
$I-B_U(I+sN)$ is affine in the indeterminate $s$, say $A+sB$.
Exact Laurent multiplication gives
\begin{align*}
 A&=-\frac{(t-1)(t^2-t+1)
 (t^8-2t^7+2t^6+t^4+2t^2-2t+1)}{t^2},\\
 B&=tP^*(t),\qquad t^2\big((1+t)A+B\big)=P(t).
\end{align*}
These identities are also verified as polynomial identities in the
supplement, before any parameter is substituted. For $k=3+6n$,
$S_k=(1+t^k)/(1+t)$, and therefore
\[
 t^2(A+S_kB)=\frac{P(t)+t^{k+3}P^*(t)}{1+t}.
\]

For completeness, this cofactor computes the knot Alexander polynomial
up to a Laurent unit. The Fox Jacobian of the Artin map $H$ in
\eqref{eq:hurwitz}, after abelianizing every meridian to $t$, is the
displayed $G_i$ block. Chronological application of the literal maps
$T$ gives the Jacobian $B^{-1}$, where $B=B_UG_2^k$.
The closed-braid presentation identifies the four final generators
with the four initial ones. Its product-preserving Artin action makes
one of these four relations redundant. Deleting that relation and one
generator column gives an Alexander cofactor of $I-B^{-1}$.
Since $\det B$ is a Laurent unit and
\[
 \operatorname{adj}(I-B^{-1})=-\det(B)^{-1}\operatorname{adj}(I-B)
\]
for these four by four matrices, the cofactor of $I-B$ computes the
same polynomial up to a Laurent unit. The adjugate identity follows
from $I-B^{-1}=-B^{-1}(I-B)$ and
$\operatorname{adj}(I-B)B=\operatorname{adj}(I-B)$.
The latter holds because $I-B$ is singular: each $G_i$ fixes the all-one
column. This also explains the convention choice without changing the
knot.

The numerator of \eqref{eq:alexander} vanishes at $-1$, since
$P(-1)=-61$ and $P^*(-1)=61$. It is reciprocal of degree $15+6n$,
with constant and highest coefficient one, so division by $1+t$ gives
a reciprocal polynomial of degree $14+6n$ with the stated ends.
Its value at one is $P(1)=1$. Differentiating the numerator at $-1$
gives $\Delta_{K_n}(-1)=-117-366n$, proving the determinant formula.
\end{proof}

The Alexander genus bound now gives $g(K_n)\ge7+3n$.
The standard braid Seifert surface has four disks and $17+6n$ bands,
so its genus is $(17+6n-4+1)/2=7+3n$, giving equality.
For any oriented diagram with $c$ crossings and $s$ Seifert circles,
its canonical surface gives $c\ge2g(K_n)+s-1$.
Yamada's theorem \cite[Theorem~3, printed p.~347]{yamada} identifies the
minimum possible $s$ with braid index, which is four here. Thus every
diagram has at least $17+6n$ crossings, and \eqref{eq:braid} attains
this bound. In particular the prime knots $K_n$ are pairwise distinct.

\section{The uniform rank-one obstruction}

Let $\rho:G_{K_n}\to GL(V)$ be any finite-dimensional representation
over a field $k$ whose meridian displacement has rank one. All meridian
matrices have a common nonzero determinant $\mu$. Write
\[
 M_i=I-u_if_i,\qquad f_i(u_i)=1-\mu,
\]
where $u_i$ and $f_i$ are nonzero. Pairs can be rescaled as
$(u_i,f_i)\mapsto(\lambda u_i,\lambda^{-1}f_i)$.
Let $r$ be the dimension of the span of all meridian roots, and $q$
the dimension of the span of all meridian covectors. Equivalently
$r=\dim\sum_g\operatorname{im}(\rho(g)-I)$ and
$q=\operatorname{codim}V^{G_{K_n}}$.

\begin{theorem}\label{thm:paired}
For every field, every $n\ge0$, and every such representation,
\[
 r\le4,\qquad q\le4,\qquad \min(r,q)\le3,\qquad r+q\le7.
\]
No irreducibility, preserved form or nonsingular evaluation matrix is
assumed, and $\mu=1$ is permitted.
\end{theorem}
\begin{proof}
The four initial meridians generate the closed-braid group, so their
root and covector spans are invariant under them and their inverses.
They contain the root and covector of every conjugate meridian;
therefore $r,q\le4$. Suppose both dimensions were four.
We will force a fixed three-dimensional contradiction.

Here is a label table for the fourteen-letter prefix. A row
$(o,a,b,\varepsilon)$ means $M_b=M_o^\varepsilon M_aM_o^{-\varepsilon}$.
The initial strand labels are $(1,0,8,13)$.
\[
\begin{array}{c|rrrr@{\qquad}c|rrrr}
0&0&1&2&-&1&8&2&3&-\\
2&3&13&14&+&3&14&3&4&+\\
4&4&8&9&-&5&9&14&15&+\\
6&15&9&10&+&7&10&4&5&-\\
8&10&0&1&-&9&1&10&11&-\\
10&5&11&12&-&11&12&5&6&-\\
12&6&12&13&-&13&13&15&16&+
\end{array}
\]
The final labels of this prefix are $(1,6,16,13)$; the outer labels are
identified by closure because the tail touches only positions two and
three. The two occurrences of label $13$, and likewise those of $1$,
express precisely these outer closure relations, independently of the
length of the tail.

Put $A=M_2$, $B=M_8$, $C=M_{13}$, $D=M_0$.
Replacing initial $M_1$ with $A=D^{-1}M_1D$ is a triangular change in
both root and covector families, with diagonal entries one. Thus the
four pairs for $A,B,C,D$ remain independent separately.
Let $H$ be the root span of $A,B,C$ and $F$ their covector span.
Write $D=(v,\delta)$, so $v\notin H$ and $\delta\notin F$.
Rows $1$ through $7$ and row $13$ generate pairs in $H,F$; in particular
this applies to $X=M_{10}$, $P=M_5$ and $Y=M_{16}$.

The first two folds are written explicitly here to record their
independence of the tail and of the parameter $1-\mu$.
Rows $0,8$ give $DAD^{-1}=X^{-1}DX$. If the pairs of $A,X$ are
$(a,f),(x,\xi)$, equality of the nonzero rank-one operators gives some
$\lambda\ne0$ with
\begin{align*}
 a-\delta(a)v&=\lambda(v+\mu^{-1}\xi(v)x),\\
 f+\mu^{-1}f(v)\delta&=\lambda^{-1}(\delta-\delta(x)\xi).
\end{align*}
Comparing components outside $H,F$ gives
$\lambda=-\delta(a)$ and $f(v)=\mu\lambda^{-1}$.
The remaining components yield
\[
 a=\lambda\mu^{-1}\xi(v)x,\quad
 f=-\lambda^{-1}\delta(x)\xi,\quad
 \delta(x)\xi(v)=-\mu.
\]
Multiplying proves $af=x\xi$, so $X=A$. The group relation becomes
$DAD=ADA$, which also gives $M_{11}=D$ from row $9$.

Rows $10,11,12$ are three maps $T$:
\[
 (D,P)\longmapsto(P,N)\longmapsto(N,Q)\longmapsto(Q,C),
 \quad N=M_{12},\quad Q=M_6.
\]
Choose the pair of $P$ as $(p,\pi)$ in $H,F$, and put
$x=\pi(v)$, $y=\delta(p)$, $h=1+xy/\mu$.
Direct rank-one conjugation gives
\[
\begin{array}{c|c|c}
 &\text{root}&\text{covector}\\\hline
N&v+(x/\mu)p&\delta-y\pi\\
Q&yv+hp&-(x/\mu)\delta+h\pi\\
C&hv+(x/\mu)(1+h)p&h\delta-y(1+h)\pi.
\end{array}
\]
The pair on $C$ differs from its seed pair by at most reciprocal
nonzero scaling. Its root lies in $H$ and its covector lies in $F$.
Either fourth component gives $h=0$. Hence $xy=-\mu$, with $x,y$
nonzero. The last two rows then give $Q=D$ and $P=C$ as matrices.
These comparisons never divide by $1-\mu$.

The prefix middle pair is now $(D,Y)$, and closure requires its tail
image to be $(D,B)$. Each of the $3+6n$ maps $T$ preserves the ordered
product. Consequently
\[
                    DY=DB,\qquad\text{so }Y=B.
\]
This cancellation replaces any argument restricted to three tail
crossings. It is valid for the entire family.

The remaining relations are exactly
\begin{equation}\label{eq:core}
\begin{array}{rrrr}
(8,2,3,-),&(3,13,14,+),&(14,3,4,+),&(4,8,9,-),\\
(9,14,15,+),&(15,9,2,+),&(2,4,13,-),&(13,15,8,+).
\end{array}
\end{equation}
Under the bijection
\[
\begin{array}{c|rrrrrrrr}
\text{here}&2&3&4&8&9&13&14&15\\\hline
\text{reduced system}&11&10&9&0&12&8&6&5
\end{array}
\]
these are the reduced $13a_{392}$ equations of
\cite[Lemma \texttt{lem:cores}, $392$ case]{extensions}, with
$0=4$ and $6=7$ there. Solving a row for the other under-arc reverses
its sign; no field parameter is inverted.
All eight core pairs lie in $H,F$, and the core contains the original
independent triples $A,B,C$. Both full core spans thus have dimension
three. The reduced system's standard triple on arcs $0,6,12$
propagates all eight pairs, so its root and covector spans are
independent separately as well. The cited paired-span lemma forbids
precisely this simultaneous dimension-three situation, over every
field and every $\mu\ne0$. This uses its $392$ case alone, including
$\mu=1$; it is stronger than just a three-dimensional irreducibility
obstruction. We obtain a contradiction and hence $\min(r,q)\le3$.
Together with $r,q\le4$ this gives $r+q\le7$.
\end{proof}

For an irreducible representation the nonzero displacement span is
all of $V$ and the common fixed subspace is zero. Thus $r=q=\dim V$,
and Theorem~\ref{thm:paired} gives dimension at most three.
Cornwell's arbitrary-field augmentation-to-irreducible-KCH
correspondence \cite[Theorem~1.2 and Definition~1.8]{cornwell} therefore
bounds every admitted augmentation rank by three over every field.
For a completely reducible representation exactly one irreducible
summand can have nonzero meridian displacement. All other summands
are trivial by meridional normal generation. In particular this
applies to finite complex images and gives $r=q\le3$ there.

\section{An attaining representation}

Set $\phi^2=\phi+1$ with $\phi=(1+\sqrt5)/2$ and put
\[
 u=(-2,0,0),\quad v=(-1,\phi-1,-\phi),\quad
 w=(\phi-1,\phi,-1).
\]
Each vector has squared Euclidean norm four. For such a vector define
$R_z=I-zz^{\mathsf T}/2$. Place
\[
                 (R_u,R_v,R_w,R_w)
\]
on the four top strands. Substitution in \eqref{eq:hurwitz} gives
exactly the same tuple after the prefix $U$.
Also $(R_vR_w)^3=I$, with $R_vR_w\ne I$; this follows either by
multiplying the three by three matrices or by $v\cdot w=2$ and the
norms four. Thus $T^3$ fixes the middle pair, and every tail
$T^{3+6n}$ fixes it. This constructs a representation of every $G_{K_n}$
with meridian parameter $-1$.

The three vectors form a basis with determinant $-4$, and their
nonorthogonality graph is connected. An invariant subspace has an
invariant orthogonal complement, because these matrices are unitary.
Each reflection root lies wholly in one of the two summands; connected
nonorthogonality forces all three basis roots into the same summand.
This proves complex irreducibility and dimension three.
Cornwell's correspondence gives $\ar(K_n,\C)\ge3$, completing the
augmentation equality.

Finiteness and the class bound can be checked without a group name.
The matrices preserve the following set of thirty vectors:
\[
 \mathcal R=\{(\pm2,0,0)\text{ and its cyclic shifts}\}
 \ \cup\ \{\text{cyclic shifts of }(\pm1,\pm(\phi-1),\pm\phi)\}.
\]
All eight choices of signs occur in the second set. The action is
faithful because $\mathcal R$ spans. Exact enumeration gives image
order $120$ and a single class of fifteen reflections, one per root
line. Its three displayed generators generate the image, whereas
every pair of the fifteen reflections generates a subgroup of order
at most ten. Thus this meridian class has generating rank three.
The finite complex domination theorem
\cite[Theorem \texttt{thm:domination}]{domination} bounds the class rank
of every finite complex rank-one-meridian image by augmentation rank.
The largest such bound for each $K_n$ is consequently exactly three.

\section{Singular Gram forms and infinite Coxeter quotients}

A restriction to nonsingular evaluation matrices would leave a gap for
canonical forms of infinite Coxeter groups. The next elementary
extension removes that restriction.

\begin{lemma}[Extending a normalized reflection system]\label{lem:lift}
Let $k$ have characteristic different from two. Suppose knot meridian
matrices are normalized symmetric-form reflections
\[
 R_a(z)=z-B(a,z)a,\qquad B(a,a)=2,
\]
for one symmetric form $B$, possibly singular. Let $S$ be their
finite-dimensional total root span, of dimension $r$. There is a
representation of the same knot group on a space of dimension $2r$
preserving a nondegenerate symmetric form, whose meridian roots have
span dimension $r$ and whose meridian covectors also have span
dimension $r$. This representation need not factor through the
original image group.
\end{lemma}
\begin{proof}
Each $R_a$ preserves $B$ and is its own inverse. Conjugation gives
$R_aR_bR_a=R_{R_ab}$. In a Wirtinger relation this equals the
reflection $R_c$. Equality of the nonzero rank-one displacements
implies $c=\lambda R_ab$. Both roots have norm two, so $\lambda^2=1$,
even when $B$ is singular. Thus every crossing relation is the root
identity $c=\pm R_ab$.

Choose a basis of $S$ and let its Gram matrix be $G$. On $S\oplus S^*$
use
\[
 \widetilde B=\begin{pmatrix}G&I\\I&0\end{pmatrix},\qquad
 \widetilde B^{-1}=\begin{pmatrix}0&I\\I&-G\end{pmatrix}.
\]
Keep all roots in the first summand. Their norms and their mutual
pairings are unchanged. Every crossing root identity therefore
remains true, and the extended reflections satisfy the full
Wirtinger presentation. Nondegeneracy makes the map from a root to
its covector injective; hence both spans have dimension $r$.
\end{proof}

\begin{corollary}\label{cor:coxeter}
For every $n\ge0$, any representation of $G_{K_n}$ by normalized
reflections in one symmetric form over a field of characteristic
different from two has total root span at most three. In particular
every surjective Coxeter reflection-class quotient of $G_{K_n}$ has
Coxeter rank and meridian-class generating rank at most three.
\end{corollary}
\begin{proof}
The four generating meridians make the root span at most four, even
if the ambient space has infinite dimension. If it were four,
Lemma~\ref{lem:lift} would give an eight-dimensional representation
with $r=q=4$, contrary to Theorem~\ref{thm:paired}.

For a Coxeter quotient use its canonical real geometric
representation, with simple roots normalized to norm two; this form
may be singular \cite{davis}. Let $S$ be the span of the roots of four
generating meridian reflections. It is stable under their reflections,
and the whole image acts trivially on the quotient by $S$. Thus every
simple reflection acts trivially on that quotient. Its displacement
image is its simple root line, so every simple root belongs to $S$.
The simple roots are a basis, proving the Coxeter rank is at most
three by the first assertion. This also covers a Coxeter system
initially allowed to have infinitely many simple generators.

The abelianization of a Coxeter group is one copy of $C_2$ per
connected component of its odd-labelled diagram. A knot group is
normally generated by one meridian, so its quotient's abelianization
is cyclic; the odd diagram is connected. For an odd edge of label $m=2k+1$, the relation $(st)^m=1$ gives
$(st)^k s(st)^{-k}=t$. Thus adjacent simple reflections along an
odd edge are conjugate, and all simple reflections are conjugate.
A reflection means a conjugate of a simple reflection. Therefore the
meridian class contains the entire simple generating set and has
generating rank at most three.
\end{proof}

The normalization in Lemma~\ref{lem:lift} is essential to its scope.
It treats canonical Coxeter reflections and normalized symmetric
reflection systems; it is not a statement about all rank-one
isometries of arbitrary degenerate forms.

\section{Adding trefoils}

\begin{corollary}\label{cor:sums}
For all integers $n,s\ge0$, let $J_{n,s}=K_n\#(\#^s3_1)$. Then
\[
 \mr(J_{n,s})=b(J_{n,s})=4+s,\qquad
 \ar(J_{n,s},\C)=3+s.
\]
Every finite complex rank-one-meridian image has meridian-class
rank at most $3+s$. Every Coxeter reflection-class quotient also has
meridian-class rank at most $3+s$.
\end{corollary}
\begin{proof}
Keep the transitive eight-sheet action \eqref{eq:a8}. Its meridian
has six cycles and defect $\delta=8-6=2$. On each trefoil factor send
all meridians to that same involution. The corresponding disconnected
branched cover has two double-cover components and four one-sheet
components. Each double cover of a trefoil has $\F_3$ first homology
of dimension one: the trefoil orbifold group
$\langle a,b\mid a^2=b^2=(ab)^3=1\rangle$ is $S_3$, with even
kernel $C_3$.
The sum sphere lifts to six spheres joining the connected seed cover
to these six components in a star tree. Mayer--Vietoris therefore
adds two to first homology per trefoil. If the seed dimension is
$h_n\ge0$, the sum cover has dimension $h_n+2s$.
The arbitrary-cycle branched-cover rank inequality
\cite[Theorem \texttt{thm:general}]{cover} gives
\[
 \mr(J_{n,s})\ge\left\lceil\frac{8-1+h_n+2s}{2}\right\rceil
 \ge4+s.
\]
Joining bridge presentations gives $b(J_{n,s})\le4+s$.

Cornwell's connected-sum theorem
\cite[Theorem~1.9]{cornwell} bounds augmentation rank above by the sum
of the two ranks minus one, and constructs an attaining representation
when the factors have attaining KCH representations with the same
meridian parameter. The preceding rank-three representation has
parameter $-1$. A rank-two trefoil representation with this parameter
is given by
\[
 A=\begin{pmatrix}-1&1\\0&1\end{pmatrix},\qquad
 B=\begin{pmatrix}1&0\\1&-1\end{pmatrix}.
\]
They satisfy $ABA=BAB$ and generate the full two by two matrix algebra,
so they are irreducible. Iteration proves augmentation rank $3+s$.
Finite complex domination then gives the stated finite-class upper
bound.

For a Coxeter quotient, the root span contributed by the $K_n$ factor
is at most three by Corollary~\ref{cor:coxeter}, applied to its
restricted normalized reflection system. A trefoil group is generated
by two meridians, one identified with a meridian of the existing sum.
It adds at most one new root direction. Thus the entire root span is
at most $3+s$. The invariant-span argument in the preceding corollary
bounds Coxeter rank by $3+s$, and the same connected odd-diagram
argument bounds the generating rank of the meridian class.
\end{proof}

\section{Exact verification and dependencies}

The finite data supplement contains the literal fourteen-letter prefix,
step-six tail rule, four permutations \eqref{eq:a8}, the three-dimensional
roots, and the coefficients of $P$. Its standard-library verifier
checks the full $A_8$ enumeration, the two exact pair periods, formal
Laurent polynomial identities before parameter substitution, prefix
signed labels, the first two fold identities and the exact reduced
core bijection. It checks the inverse of the extended Gram matrix
symbolically. Rejected corrupt controls alter a tail period, a
permutation, an Alexander coefficient, a run count, a prefix label,
and a positive root. Optional topology replay reconstructs selected
literal braid closures independently. Representative values test
implementation; the universal statements use the period induction,
affine cofactor identity, ordered-product cancellation, and written
Gram extension.

The essential imported mathematical conclusions are the $392$ paired
contradiction, the published homogeneous-braid primeness theorem,
Yamada's theorem, Cornwell's correspondence and sum theorem, the
branched-cover rank inequality, and the finite complex domination
theorem. The accompanying receipt identifies exact source editions and
separates exact replay from these written premises. The assertion about
all Coxeter targets uses the new Gram extension rather than any unit
ideal computation at the initial parameter. No formal verification,
external specialist review or worldwide-priority assertion is implied.

\paragraph{AI assistance and review.}
ChatGPT (OpenAI) assisted with mathematical arguments, literature checks,
exposition, proof review and verification code.

\begin{thebibliography}{9}\small
\bibitem{flo} P.~Feller, L.~Lewark and M.~Orbegozo Rodriguez,
\emph{Homogeneous braids are visually prime}, J.~Topology 18 (2025),
e70040, \href{https://doi.org/10.1112/topo.70040}{doi:10.1112/topo.70040}.
Retained arXiv:2408.15730v2 (30 September 2025), Theorem~1.1,
\S5.1--5.2 and Lemma~5.1(i).
\bibitem{yamada} S.~Yamada, \emph{The minimal number of Seifert circles
equals the braid index of a link}, Invent.~Math.~89 (1987), 347--356,
Theorem~3, printed p.~347.
\bibitem{cornwell} C.~R.~Cornwell, \emph{KCH representations,
augmentations, and $A$-polynomials}, J.~Symplectic Geom.~15 (2017),
983--1017. Retained arXiv:1310.7526v3, Theorems~1.2 and~1.9,
Definition~1.8 and Corollary~1.12.
\bibitem{extensions} \emph{Signed generic rank-one obstructions for
sixteen knots}, companion \path{GENERIC_KCH_EXTENSIONS.tex},
Lemma \texttt{lem:cores}, $392$ subcase only.
\bibitem{domination} \emph{Finite complex reflection-class bounds are
dominated by augmentation rank}, companion
\path{FINITE_COMPLEX_AUGMENTATION_DOMINATION.tex},
Theorem \texttt{thm:domination}.
\bibitem{cover} \emph{Meridional rank bounds from arbitrary-cycle
branched covers}, companion
\path{GENERAL_BRANCHED_COVER_RANK_BOUNDS.tex},
Theorem \texttt{thm:general}.
\bibitem{davis} M.~W.~Davis, \emph{The Geometry and Topology of
Coxeter Groups}, second edition, 2025. Canonical geometric
representation: \S6.12, equations (6.32)--(6.33), Lemma~6.12.3
and Corollary~6.12.4, printed pp.~108--109 (PDF pp.~127--128).
Only existence and the simple-root basis are used, not faithfulness.
\end{thebibliography}
\end{document}
